\documentclass[pmlr]{jmlr}% new name PMLR (Proceedings of Machine Learning Research)
% Template adapted for the Workshop on Emerging AI Technologies for Music, as part of AAAI
% https://amaai-lab.github.io/EAIM2027/

 % The following packages will be automatically loaded:
 % amsmath, amssymb, natbib, graphicx, url, algorithm2e

 %\usepackage{rotating}% for sideways figures and tables
\usepackage{longtable}% for long tables

 % The booktabs package is used by this sample document
 % (it provides \toprule, \midrule and \bottomrule).
 % Remove the next line if you don't require it.
\usepackage{booktabs}
 % The siunitx package is used by this sample document
 % to align numbers in a column by their decimal point.
 % Remove the next line if you don't require it.
\usepackage{siunitx} % newer version % newer version
 %\usepackage{siunitx}

% my code listings
\usepackage{listings}
\lstset{language=lisp,
 columns=flexible,
 breaklines=true,
 breakatwhitespace=true,
 identifierstyle=\ttfamily,
 keywordstyle=\ttfamily,
 basicstyle=\footnotesize\ttfamily,
 showstringspaces=false,
}

 % The following command is just for this sample document:
\newcommand{\cs}[1]{\texttt{\char`\\#1}}

 % Define an unnumbered theorem just for this sample document:
\theorembodyfont{\upshape}
\theoremheaderfont{\scshape}
\theorempostheader{:}
\theoremsep{\newline}
\newtheorem*{note}{Note}

 % change the arguments, as appropriate, in the following:
\jmlrvolume{YYY} % Change this in Camera Ready Version
\jmlryear{2027}
\jmlrworkshop{EAIM2027 at AAAI}

\title[Short Title]{Template for EAIM2027\titletag{\thanks{sample footnote}}}

 % Use \Name{Author Name} to specify the name.

 % Spaces are used to separate forenames from the surname so that
 % the surnames can be picked up for the page header and copyright footer.
 
 % If the surname contains spaces, enclose the surname
 % in braces, e.g. \Name{John {Smith Jones}} similarly
 % if the name has a "von" part, e.g \Name{Jane {de Winter}}.
 % If the first letter in the forenames is a diacritic
 % enclose the diacritic in braces, e.g. \Name{{\'E}louise Smith}

 % *** Make sure there's no spurious space before \nametag ***

% Author information should be kept anonymous for the double-blind policy. Add in this information for the camera ready version.
 % Two authors with the same address
  \author{\Name{Anonymous 1\nametag{\thanks{with a note}}} \Email{abc@sample.com}\and
   \Name{Anonymous 2} \Email{xyz@sample.com}\\
   \addr Address}

 % Three or more authors with the same address:
 % \author{\Name{Author Name1} \Email{an1@sample.com}\\
 %  \Name{Author Name2} \Email{an2@sample.com}\\
 %  \Name{Author Name3} \Email{an3@sample.com}\\
 %  \Name{Author Name4} \Email{an4@sample.com}\\
 %  \Name{Author Name5} \Email{an5@sample.com}\\
 %  \Name{Author Name6} \Email{an6@sample.com}\\
 %  \Name{Author Name7} \Email{an7@sample.com}\\
 %  \Name{Author Name8} \Email{an8@sample.com}\\
 %  \Name{Author Name9} \Email{an9@sample.com}\\
 %  \Name{Author Name10} \Email{an10@sample.com}\\
 %  \Name{Author Name11} \Email{an11@sample.com}\\
 %  \Name{Author Name12} \Email{an12@sample.com}\\
 %  \Name{Author Name13} \Email{an13@sample.com}\\
 %  \Name{Author Name14} \Email{an14@sample.com}\\
 %  \addr Address}


 % Authors with different addresses:
 % \author{\Name{Author Name1} \Email{abc@sample.com}\\
 % \addr Address 1
 % \AND
 % \Name{Author Name2} \Email{xyz@sample.com}\\
 % \addr Address 2
 %}

\editors{K. Bhandari, S. Zheng, J. Shier, D. Herremans, S. Colton}

\begin{document}

\maketitle

\begin{abstract}
This is the abstract for this article.
\end{abstract}
\begin{keywords}
List of keywords
\end{keywords}

\section{Introduction}
\label{sec:intro}


In the 2020's, machine learning, and in particular deep learning, has inarguably become the dominant paradigm
for artificial intelligence implementations, including in the field of music. 
However, classical AI techniques, such as expert systems implemented using constraint satisfaction solvers, have 
characteristics that make them that make them well suited to a variety of musical problems,
and on modern machines, can be used in soft real time contexts where previously this would have been impractical.

Constraint-based systems have the particular advantage that, while the computation time necessary
to find a solution may be indeterminate, the solution itself (should one be found) can be guaranteed to meet formally
specified criteria. For the music teacher or algorithmic composer, this can be a critical feature.
For example, in the creation of musical training tools, a tool author may wish to specify that the 
symbolic musical content (perhaps presented as a score or played aloud using standard synthesizers)
follow specific voice leading rules given a particular chord progression. 
Or a jazz ear training exercise may use rules to specify that a walking bassline will only use
specific generative rules, thus ensuring there is no ambiguity as to the underlying chord progression.
In the case of a pedagogical exercise builder, it is not sufficient that the material generated be almost
guaranteed to follow the rules specified as false-positive or false-negatives will be extremely frustrating
to the student. 

Or finally, a composer may wish to have stochastically generated material follow specific rules, with
the option of interactively changing these rules as a piece is being played.

We present here S4M-CSP, an implementation of a domain specific constraint solver and accompanying symbolic music theory
library, implmented in s7 Scheme Lisp CITE BILL and run in the Max/MSP musical programming environment using the 
open-source Scheme for Max extension. The authors are using this implementation to create both training
exercises and algorithmically generated material for interactive electro-acoustic music.

Of note is that the solver is implemented entirely in high-level Scheme, using a brute-force
backtracking search. While this is computationally inefficient compared to the possibilities of optimized
contraint solvers CITE SOMETHING, this has the advantage that constraints themselves can be
expressed as simple and highly-readable Scheme functions, providing the advantages of constraint-based 
systems to users with no specialized mathematical, and limited computer science, background. 
As these are specified in uncompiled text code and run directly in the Scheme for Max interpreter in 
the Max interactive environment, the (brave) user can even update constraints and
domains dynamically in real time. 

Additionally, the implmentation described includes a recursively-generated record of all paths
previously taken, enabling the system to use randomly (or weighted randomly) chosen paths, producing
a different solution on each run, where the potential solutions are many.

The results suggest that on modern computers, the use of even simplistic constraint solvers implemented
in user-friendly formats can be a usable and user-friendly way of solving rules-based symbolic musical
problems for both pedagogical and compositional purposes, and that this can provide valuable tools for
pedagogues, composers, and performers. 

\section{Related Work }
\label{sec:intro}

The use of constraint programming for solving music composition problems (i.e., computationally generating 
material that obeys defined rules) has an extensive and long history, dating as far back as the 1950's with
Hiller and Isaacson's landmark composition, Illiac Suite (1957). The compositional approach
to the Illiac Suite was one of "generate and test", in which separate algorithms were used
to generate material and to filter it for acceptability by testing against programmatic constraints.
(Nierhaus, 2009)

Gecode, one of the major general constraint satisfaction programming platforms, has been in active development
since the early 2000s and is used in musical contexts such as constraint generated harmonic progressions,
for example in Sprockeels' "Harmonizer Project" (Sprockeels, 2025)

More related to the work in question, Lisp as a language has a rich history in musical constraint programming. 
For example, the PMC library, part of the PWGL computer-assisted-composition
framework  allows composers to define constraints in Common Lisp. (Sandred, 2010).
This solver is
is also available in the OpenMusic platform as the library OMCS, and, like S4M-CPS,
finds solution through a backtracking search. (Anders, 2018)
Gelisp is similarly a system for defining constraints in Lisp. Gelisp is implemented as a
a Lisp wrapper around the C++ Gecode library, and also provides integrations for the OpenMusic
graphical composition platform (Toro, 2015)

Musical constraint systems are also used outside of computer-assisted-composition in various
ways, including computational musicology (the study of music using computers).
In his writings on his ground-breaking algorithmic composition system EMI (Experiments in Musical Intelligence)
David Cope discusses programmatically constraining material chosen randomnly form corpus databases
to discover the defining features of musical style.
Duane and Pardo (Duane, 2009) describe using constraint optimization techniques as part 
of computational system for auditory stream analysis from MIDI data of Classical, Baroque,
and Romantic music. 
In his widely acclaimed theory text, "The Geometry of Music", Dmitri Tymoczko describes using 
constrained, randomly generated melodic and harmonic pitches
to analyse which specific musical rules or tendencies contribute to a listener's perception 
of style and consistency.  (Tymozcko, 2011).

In music education, constraint programming techniques have been used to generate exercices
that are guaranteed to follow particular rules or stay within particular bounds, or to
create sequences of exercices at appropriate difficulty levels for optimal learning.
Examples include generation of rhythmic reading exercises (Percival, 2008), (Kuuskankare, 2013)
intonation exercises (Percival, theses), and curriculum planning (holland, 2013). 

The use of constraints for computer assisted composition is supported in the Max environment through the BACH
library. ( BACH provides a suite graphical programming objects for algorithmic composition,
inspired in part by OpenMusic.
This includes support for manipulating Lisp-like data structures (lllls, or "Lisp-like linked lists).
(Agostini, 2012)
However, the BACH library does not provide any built in Lisp interpreter, thus while users can 
use constraints in ways similar to a Lisp environment, they cannot code in a Lisp directly
to make new constraints.


\section{Language and Platform}

\subsection{Max}
S4M-CSP runs on the Max audio/visual low-code programming environment, which was first released over
30 years ago and is used widely by musicians, composers, programmers, and audio visual artists to create
music and multimedia works and interactive environments. 
Max is predominantly a visual programming environment, in which interactive programs (called "patches" in Max parlance)
are placing visual objects on a canvas and connecting them with virtual wires that pass symbolic messages or blocks 
of audio samples. However, in addition to the visual paradigm, Max supports textual programming predominantly through an integration
with JavaScript (i.e., the js object), as well as through code editing for the Gen audio objects, 
external processes through Node4Max, and Java code. 
In addition to the built in facilities provided with Max, users may use "externals", third-party objects
created using the Max SDK and built in C or C++. (CITE Puckette or Zicarelli)
Various of these also support working in code in Max through
interfaces with other languages or embedded interpreters. 


\subsection{Scheme for Max and Lisp}
S4M-CSP is a code library for Scheme for Max (S4M), an open-source extension to Max that enables users to
program Max using s7 Scheme, a Lisp dialect authored by Bill Schoetstadt of CCRMA,  and also used
in Common Music 3 and the snd audio editor. While the facilities of S4M-CSP could certainly be implemented
in the built-in JavaScript facilities, S4M was used for several reasons. 

First, as a Lisp, S4M provides first-class support for symbols as a primitive data type, where such symbols
can be used identically in both Max messages (the graphic boxes in Max for holding and manipulating symbolic sequences),
and in Scheme code. Lisp symbols are ideal for manipulating music symbols as they can be compared directly,
rather than requiring string searching and matching. (Both s7 and Max maintain a symbol table for fast 
lookup and comparison, and symbols may be used as hash-table keys for caching.) In the context of a constraint
solver, where symbolic predicates may be evaluated hundreds-of-thousands of times per solution pass,
this balance of user convenience and processing speed is a signficant advantage.

Secondly, unlike the built-in JS object, S4M can be run in either the high- or low-priority Max thread.
Max provides a high-priority thread (the "scheduler" thread) for handling audio, MIDI messages, and other
timing critical events, as well as a low-priority thread (the "main" or "GUI" thread) for display, file I/O,
and on screen events (e.g., clicking and dragging elements). The built in JS object in Max runs only in the
lower priority thread, making it inappropriate for directly scheduling events, such as one would do 
when building a sequencer in code. Because S4M can run in both threads, and because Max message syntax
is Scheme-compatible (very few syntax collisions), using S4M allows one to create a patch where the
CSP solver runs in the low thread, but as input sent to it from a high-priority S4M instance, and sends
results to the same or a second high-priority S4M instance for controlling playback. 

Finally, while the initial version of S4M-CSP is a stand-alone library, there is a rich body of related
work available to users in Lisp dialects, including in s7 itself (e.g., through Common Music 3 code).
As previously mentioned, there is a significant history of Lisp code in both algorithmic composition
and computational musicology. While s7 is a Lisp-1, and thus differs from Common Lisp in a variety of ways,
the author of s7 has made compatability with much of Common Lisp a high priority, including full support
for Common Lisp facilites such as defmacro, reader macros, keywords, hash-tables, format, and more.
Thus, while the user need not know or install anything beyond S4M and the S4M-CSP library to use the
solver, the use of S4M provides users the potential to integrate their work with numerous other Lisp 
and Lisp related musical libraries and tools, such as Common Music 3, Slippery Chicken, Incudine, 
and many more.


\section{Positioning and Contribution}
Given the breadth of CSP options for music, one might well ask why one might want to use S4M-CSP
and what novel contributions the project makes relative to the other available options
for constraint programming for music.

\subsection{Platform Accessibility}
First, S4M-CSP is intended to be accessible; it aims to provide the power and flexibility of
constraing programming to the widest possible range of computer music composer-programmers and
music educators. Max is arguably the most widely deployed computer music programming platform, 
by a significant margin, as it is not only used worldwide for audio, music, and visual
programming, but is shipped with the hugely successful Abelton Live digital audio workstation
as Max for Live. While users need to install the S4M extension to Max, this is a well-documented
and trivial step. The authors believe this is significantly more accessible to the average user
than the competing offerings (e.g., OpenMusic, Patchworks, etc.) which require the user to 
adopt a specific environment rather than being able to run a solver within an already mainstream
environment.

\subsection{Ease of Use}
Second, using S4M-CSP requires less computer science background than many other options. 
Programming constraints in s7 Scheme with the provided music theory library is very straightforward.
Scheme is a very minimal language, and as a Lisp 1, it has a simpler syntax for higher order 
programming than Common Lisp (i.e., there is no second namespace for functions from variables).
The Scheme for Max project provides beginner tutorials for Scheme and free and open-source
resources for learning Scheme abound. Additionally, interfacing with the solver from Max
is simple as one can generate Scheme expressions directly in Max messages, simplifying the 
process of creating real-time interfaces to the project.

TODO: image of s-exp in a Max patch.

TODO: code of generating a constraint

\subsection{Exensibility}
Third, both S4M and s7 Scheme are open-source and implemented in ANSI-C.
Thus, the intrepid user can extend the solver at the language level, implementing
new s7 predicates in C where performance optimizations warrant this. 
In fact, one of the major reasons s7 was chosen by the author of Scheme for Max 
was its simple foreign function interface and the ease with which it can be embedded in C projects.
(CITE S4M Website)

\subsection{Accurate Music Symbol Support}
Finally, the S4M-CSP project and its underlying theory library are designed to support
working in music notation where accuracy of symbol choice is important.
To this end, S4M-CSP makes no assumptions as to when enharmonic equivalence is considered
good enough. That is to say, the theory library understands that an augment-fourth
above a C# is an F##, and that a solution finding a G should not necessarily be
considered good. While this is not an important consideration in contexts 
where a user "downsamples" to MIDI note numbers, in a situation where the project
is used to make training exercises (e.g., harmony, aural skills), this distinction
is critical and libraries that do not make it are unfit for purpose.

The underlying theory library is implemented such that it correctly determines
intervals regardless of the number of double sharps and flats, and uses intervals
up to double-diminished and double-augmented correctly. While this is computationally
much more expensive than more naive matching, the library uses aggressive caching
throughout to negate that disadvantage. Music interval calculations are, by definition,
pure functions; a double-augmented fourth above F is always B#. Thus, trading
space for time makes sense on modern computers, and the S4M-CSP engine does this
for interval, chord, and membership functions (i.e. checking whether the third of
a G major 7 chord is B should always return true).


\section{Objective}
\label{sec:intro}

The high-level goal of S4M-CSP is to enable users of the Max audio/visual 
programming environment to generate symbolic pitch sequences (i.e., the pitch component of melodies) 
given a symbolic chord progression and an arbitrary set
of rules relating the generated symbolic pitches with both the chord progression data and the neighbouring
symbolic pitches. 

Of note is that the representation of pitches, relationships, and chords is entirely as Lisp symbols,
such as C\#4, IImin7, dim-5, etc. While MIDI note numbers may be also used (by converting symbols to MIDI note numbers),
the project includes a music theory library enabling easy symbolic manipulation in terms that are 
common in Western music theory.  Thus the musically educated user must only know the symbolic 
conventions used, rather than any details of the underlying implementation. The underlying functions
accept both pitch classes and pitches, and are implemented to use correct intervallic spelling, including
arbitrary numbers of sharps or flats. 
For example, (prf-4 'C4) returns 'F4, (aug-4 'C\#) returns 'F\#\#, (aug-4 'C\#\#4) returns 'C\#\#\#5.
In order to make building systems with Max possible, all spellings use standard ASCII 
(i.e., no special unicode symbols are used for accidentals).

\begin{itemize}
  \item Upper case letters specify a pitch letter in letters space
  \item Upper case Roman numerals incidate functional roots
  \item Lower case b and the \# sign indicate accidentals
  \item double sharps or flats used multiple sequential 'b' and '#' characters.
\end{itemize}

Rules predicates are specified as Scheme callables (functions or macros) that have a set
signature providing access to the solver's context variables and to the solver object
itself. The music theory library provides lower level functions for making predicates 
out of intervals, chord spellings, and Roman numeral realizations. 

\subsection{Test Application}
The test application for S4M-CSP is a program that builds idiomatic jazz bass lines
as lists of pitch class symbols. The solver expects the following inputs (i.e., context variables)

\begin{itemize}
  \item A Roman numeral progression such as '(IImin7 V7 Imaj7), '(IV7 V7 I7), etc.
  \item A key in which the progression is rendered such 'F, 'G, etc.
  \item A domain of possible solution values: '(C1 'C#1 'Db1 ...)
  \item A list of constraint predicates as Scheme functions
\end{itemize}

The solver will attempt all combinatorial paths through the set of domain values.

\subsection{Specific Requirements}
Specific requirements of the system and of the test application are the following:

\begin{itemize}
  \item Constraints are specified in Scheme as boolean predicate functions.
  \item Constraints may be either local (constraining a single pitch solution against other context variables)
     or global (constraining between solution pitches, such as all-diff). 
  \item Multiple runs of the solver will create different solution sequences, i.e., the is not looking for the first or best solution, but rather a stochastically chosen acceptable solution.
  \item The solver will find a solution, if one exists, or return false if no solution is possible.
  \item Solutions are returned as a Lisp list of pitch-class symbols, such as '(C4 Eb4 Bb3 G3)
  \item The domain over which the solver runs may be set dynamically prior to each pass, allowing the
        user to constrain the possible pitch values by both range and key
\end{itemize}

Additionally, a goal of the project is to explore whether, on modern computers, this approach
to generating pitch sequences can be practically used in a soft-real time context, such as
generating an idiomatic bass line that plays forever. To this end,
the solver accepts a starting and ending pitch class for each round of solving, used 
for targeting a note, and the underlying music theory library uses hash-table caches
of every harmonic calculation. 


\section{Implementation}

Fill me given how much space we have after all the other shit

\section{Limitations}


\subsection{Propagation}

The current implementation is a backtracking brute-force solver and does not
include any forward propagation of constraints. 
When constraints fail at a given depth, the solver will next try alternative domain values,
and if all fail, it will backtrack up a depth. There is no implementation of pruning domain values
during a solve or of undoing a prune when backtracking up a level. 

In order to track the solve progress, allowing influencing the next attempted path,
the implementation includes a tree representation of all paths that have been explored,
with double-linked nodes such that paths from a node can be traversed up or down the tree. 
This structure could be also be used to store domain changes from visiting and assigning a node, 
which could be used to undo domain pruning when backtracking. 
However, at present the nodes-visited tree is used solely to track which paths have
been taken so that random choices can be made at each node without revisiting previous paths.

TODO: do we show the tree walker?

\subsection{Variable Selection}
The implementation of choosing the next variable to be solved (i.e., the next pitch in the solution
sequence) is naive;
it simply iterates through the note list in order. Allowing solving order
to be indicated or optimized could potentially reduce the number of paths explored as well
as provide for potentially interesting compositional algorithms. For example, in 
bass line implementation, one could work either forward or backward from the target.

\subsection{Value Selection}
Similar to variable selection, value selection at present simply chooses random potential values
from the domain list as it is set at solver initialization time.




\section{Future Work}
There are a number of areas of future work I intend to explore.

\subsection{Stochastic Heuristics}
When choosing the next candidate domain value, the current implementation either
takes the next one from the list of domain values or a random value. The addition of preference constraints
to help choose candidates in more musically meaningful ways would be a fruitful 
area to explore. This could potentially allow one to implement preference rules
such as *"if the last three notes have ascended, increase the probability of 
the next note descending"*. These rules could perhaps be implemented using
a higher-order Markov chain to enable some higher-level musical variety.

\subsection{Advanced Pruning}
Given the domain-specific constraint implementation, it would be reasonably straightforward
to add a special type of constraint intended to prune variable domains prior to solving,
based on musical rules or heuristics, which could speed up execution significantly.
For example, a user might specify that all notes up to the final target note 
should be diatonic to the tonic key. Such a constraint could
thus prune the domain for notes 1 to 3 to be diatonic only, cutting the number of 
branches to be explored down significantly.

\subsection{Constraint Propagation}
Constraint propagation will significantly speed up execution, though execution time
is already fast enough to solve lines of quarter notes as they play given a domain space
of two or three octave. This could be an issue, however, for solving longer lines with a wider 
range.

\subsection{Assignment}
Notes are assigned one at a time, but this is not an ideal heuristic for simulating
how musicians think. For example, when playing a bass line, a musician is likely
to think, *"I can walk stepwise or arpeggiate"*, determining which based on available range
and what they have recently done and evaluating whether several candidates will work 
at one time. The solver could draw from a bank of patterns to try various compound
selections first rather than brute force traversal through all paths of the state space.

\subsection{Constraint Macro Framework}
While implementing the constraints as boolean predicates is straightforward, it does
require the developer to be somewhat familiar with the implementation details of
the solver. A higher-level framework built on Lisp macros could be implemented
to reduce the amount of code necessary and to reduce how much must be understood
by the users making constraints.

\subsection{Solver Instantiation}
Creating, initializing, and setting up a solver with constraints and context variables
is currently a fairly verbose affair. This is an area where Lisp macros could
again allow a much more succint and readable API for users of the solver.

\section{Conclusion}
As a first attempt at solving musical line building as a constraint satisfaction problem
with a custom solver,
I feel the project has been a success. Execution is fast enough to allow
generating a bass line in perpetuity. While still somewhat verbose, specifying
rules in a similar way to how one thinks of musical constraints works well. 
For example, one can read the set of constraint registrations to make an
arpeggiated bass line and understand right away what the outcome is intended to be.

Breaking up the constraints into (solved) notes and (pre-assigned) context variables is 
also successful and makes domain specific problems much easier to reason about (and
to design constraints for) than would be the case with purely numerical variables that
make no destinction between musical context and the line. It is also more reflective of
how improvisers and composers think of the problem in homophonic improvised music: 
the chord progression and key are predetermined, while the notes to realize it are improvised,
choosing from harmonic constraints that determine the genre of music. For example,
a style-imitation program designed this way would use different constraints for
different eras of jazz music as the melodic and harmonic vocabulary changed over time.

The project runs well in Scheme for Max and, I believe, has the potential to be
a valuable contribution to the available offerings for computer-assisted composition 
in the Max environment, such as the Bach project. I intend to continue to work
on the project, with the intent of publishing more complete versions as part
of the Scheme for Max project.




  





\section{Citations and Bibliography}
\label{sec:cite}

The \textsf{jmlr} class automatically loads \textsf{natbib}.
This sample file has the citations defined in the accompanying
BibTeX file \texttt{eaim.bib}. For a parenthetical
citation use \verb|\citep|. For example
\citep{guyon-elisseeff-03}. For a textual citation use
\verb|\citet|. For example \citet{guyon2007causalreport}.
Both commands may take a comma-separated list, for example
\citet{guyon-elisseeff-03,guyon2007causalreport}.

These commands have optional arguments and have a starred
version. See the \textsf{natbib} documentation for further
details.\footnote{Either \texttt{texdoc natbib} or
\url{http://www.ctan.org/pkg/natbib}}

The bibliography is displayed using \verb|\bibliography|.

\acks{Acknowledgements go here.}

\bibliography{eaim}

\appendix

\section{First Appendix}\label{apd:first}

This is the first appendix.

\section{Second Appendix}\label{apd:second}

This is the second appendix.

\end{document}
